\documentclass[12pt]{article}
\usepackage[T1]{fontenc}
\usepackage[utf8]{inputenc}
\usepackage{lmodern}
\usepackage{textcomp}
\usepackage{enumitem}
\usepackage{geometry}
\usepackage{graphicx}
\usepackage{amsmath, amssymb}
\geometry{margin=1in}
\begin{document}
\begin{center}
Mathematics - K-12 Level Assessment 24 \
Total Marks: 40 \
Answer all questions.
\end{center}

\section*{Multiple Choice (10 marks)}
\begin{enumerate}[label=\arabic*.]
\item How many distinct diagonals of a convex hexagon can be drawn?
\begin{enumerate}[label=(\alph*)]
    \item 6
    \item 8
    \item 9
    \item 36
\end{enumerate}
\item Find the value of the expression 14 + (-71).
\begin{enumerate}[label=(\alph*)]
    \item -85
    \item 57
    \item 85
    \item -57
\end{enumerate}
\item The original price of a new bicycle is \$138.00. If the bicycle is marked down 15\%, what is the new price?
\begin{enumerate}[label=(\alph*)]
    \item \$20.70
    \item \$117.30
    \item \$123.00
    \item \$153.00
\end{enumerate}
\item Suppose $a$, $b,$ and $c$ are positive numbers satisfying $a^2/b = 1, b^2/c = 2, c^2/a = 3$. Find $a$.
\begin{enumerate}[label=(\alph*)]
    \item $12^(1$/7)
    \item $7^(1$/12)
    \item 1
    \item 6
\end{enumerate}
\item Roya paid \$48 for 12 cartons of orange juice. What is the unit rate per carton of orange juice that Roya paid?
\begin{enumerate}[label=(\alph*)]
    \item \$3
    \item \$4
    \item \$6
    \item \$12
\end{enumerate}
\end{enumerate}

\section*{True / False (10 marks)}
\textit{Answer True or False}
\begin{enumerate}[label=\arabic*.]
\setcounter{enumi}{5}
\item True or False: The number 300 is a powerful integer since it contains prime factors with their squares as factors.
\item True or False: The maximum value of the function f(x) = xe^\{-x\} is e.
\item True or False: The quotient of the expression 2,314 divided by 4 is 578 with a remainder of 2.
\item True or False: The solution to the equation -47 = g + 24 is g = -23.
\item True or False: The value of the expression 28 x 42 is 1,176.
\end{enumerate}

\section*{Long Form Response (20 marks)}
\textit{Answer each question with a clear and complete explanation.}
\begin{enumerate}[label=\arabic*.]
\setcounter{enumi}{10}
\item Discuss how many distinct arithmetic sequences of consecutive odd integers can be formed that sum to 240. Provide a detailed explanation of your reasoning and calculations.
\item Calculate the number of different committees of 5 people that can be formed from a group of 15. Discuss the mathematical principles involved in your solution.
\end{enumerate}

\vfill
\noindent\textit{End of Paper}
\end{document}